% =====================================================================
%  chunks/erratatable.tex —— 附录 C 原书排印问题一览
%  说明：汇总正文各页译注中指出的原书（英译本）排印问题；正文一律照原书排印、未径改。
%  注意：不能用 longtable——ctexbook 下跨页 longtable 会报 infinite glue；
%        此处沿用附录 B 术语对照表的 multicols + 悬挂缩进排法。
% =====================================================================

\chapter{原书排印问题一览}

\noindent 本表汇总全书正文译注中指出的英译本排印问题（脱字、错下标、错编号、人名拼写、数表与正文矛盾等），按原书页码排列，每条先照录原书文字（短引），再给出应为的写法或判断依据。正文一律照原书排印、未径改：凡本表所列处，读者按「说明」一栏理解即可，不必以为译文有误。少数条目并非笔误而是原书的记法习惯（如引号式交叉引用、把 $S$ 的维数写成下标），一并列出以便对照。

\small
\begin{multicols}{2}
\interlinepenalty=10000
\raggedright

\noindent\hangindent=2.2em\hangafter=1\textbf{p.21}\quad 英文 「i.e a space all of those subsets are open」有脱漏，应为 「i.e.\ a space all of whose subsets are open」（覆叠的纤维取离散拓扑）。\par

\noindent\hangindent=2.2em\hangafter=1\textbf{p.24}\quad $S^n$ 的定义印作 $\sum_{\alpha=0}^{n}(x^\alpha)^2\le1$；$\le1$ 给出的是球体，球面应为 $=1$。\par

\noindent\hangindent=2.2em\hangafter=1\textbf{p.29}\quad 相对同伦群的附贴条件印作 $f(S^n)\subset A$；$f$ 定义域中的球面部分是 $S^{n-1}$，应为 $f(S^{n-1})\subset A$。\par

\noindent\hangindent=2.2em\hangafter=1\textbf{p.41}\quad 同一句先把 $n$ 维单形记作 $\sigma_n$，随即又写成 $\sigma^n$，两种写法混用；下文一律用 $\sigma^n$（上标记维数）。\par

\noindent\hangindent=2.2em\hangafter=1\textbf{p.42}\quad 「subdivisions mu.st be such as」中的 「mu.st」是 must 的排印脱字。\par

\noindent\hangindent=2.2em\hangafter=1\textbf{p.54}\quad 卡积处印作 $a\cap b$ 而 $b$ 是上链；按数积式 (2.27) 应为 $a\cap c$。\par

\noindent\hangindent=2.2em\hangafter=1\textbf{p.59}\quad 引「式 (3.5)」，但此处讲的是上同调三元组的正合列，对应的是式 (3.6)（(3.5) 为同调版）。\par

\noindent\hangindent=2.2em\hangafter=1\textbf{p.63}\quad (ii) 中 1-胞腔列为 $\kappa_1^1,\ldots,\kappa_{2g}^1$，而 $\partial\kappa^2=2\kappa_0^1$ 里的 $\kappa_0^1$ 不在其列；对照 (iii) 应为 $2\kappa_1^1$。\par

\noindent\hangindent=2.2em\hangafter=1\textbf{p.70}\quad 例中印 $\cF_3=\cF_1/\cF_2$；由 $\cF_1\subset\cF_2$ 及下文按 $\Z/p^2/\Z/p$ 的计算可知应为 $\cF_2/\cF_1$。\par

\noindent\hangindent=2.2em\hangafter=1\textbf{p.87}\quad 与 p.88 各有一个「例 1」，编号重复；p.87 的例 1 给式 (6.4)，p.88 的例 1 才是 Hopf 分类问题的那个。\par

\noindent\hangindent=2.2em\hangafter=1\textbf{p.101}\quad (6.19) 印作 $\Ext^{1,2^i+2^j}$；由 $h_i\in\Ext^{1,2^i}$ 相乘得 $h_ih_j\in\Ext^{2,2^i+2^j}$，紧接的句子又称之为「2 线」，故上标 $1$ 应为 $2$。\par

\noindent\hangindent=2.2em\hangafter=1\textbf{p.127}\quad 式 (9.3) 首行第二个因子印作 $\Omega_G^{m_1}(X)$；由下行的反交换式 $z_1z_2=(-1)^{m_1m_2}z_2z_1$ 与乘积次数 $m_1+m_2$ 可知应为 $\Omega_G^{m_2}(X)$。\par

\noindent\hangindent=2.2em\hangafter=1\textbf{p.133}\quad 印作 $\Lambda\subset X^*$；本段讨论的是 $S$ 的自然对偶 Hopf 代数 $S^*$，且 $\Lambda\otimes\Q=S^*\otimes\Q$ 表明左边也应是 $S^*$，$X^*$ 在全书记号中并无定义。\par

\noindent\hangindent=2.2em\hangafter=1\textbf{p.169}\quad 讨论 Möbius 带处引「见图 4.6」，而图 4.6 画的是光滑单形边缘的示意，应为图 4.4。\par

\noindent\hangindent=2.2em\hangafter=1\textbf{p.184 两处}\quad (1) 式 (2.27) 的求和指标印作 $p+q=n$；$\Lambda^k$ 的分解按 $k$ 分次，应为 $p+q=k$（下文 (2.32) 即用 $p+q=k$）。(2) 式 (2.31) 第二式印作 $\Delta_c=\Delta_c$，左右端同一记号，系漏印横线；按 Kähler 恒等式应为 $\bar\Delta_c=\bar\delta_c\bar d_c+\bar d_c\bar\delta_c$，且 $\Delta_c$ 与 $\bar\Delta_c$ 相等。\par

\noindent\hangindent=2.2em\hangafter=1\textbf{p.193}\quad 式 (2.49) 的两处二项式系数均印作 $\binom{p}{p+q}$；$p,q\ge1$ 时其为 $0$，会使乘法恒为零，与 $\Omega S^n$ 上同调作为除幂代数的结构矛盾，按对称性应为 $\binom{p+q}{p}$。\par

\noindent\hangindent=2.2em\hangafter=1\textbf{p.200}\quad 式 (2.60) 右端印作 $\hat l^E(N^n_{(0,\dots,0)})=0$，与左端已严格大于 $0$ 矛盾；按上下文应为 $\hat l^E(N^n_{(1,0,\dots,0)})=\kappa_1$。同段把同伦的靶空间印作 $F$，应为覆叠空间 $\tilde M$。\par

\noindent\hangindent=2.2em\hangafter=1\textbf{p.201}\quad 脚注 5 把该泛函印作 「Chern-Simon functional」（Simon 无 s），通行写法为 Chern--Simons。\par

\noindent\hangindent=2.2em\hangafter=1\textbf{p.202}\quad (2.66) 前一行正文印 $C^q(A)$（上标），而 (2.66) 中印 $C_q(A)$（下标）；按本节记号与 (2.65)，正确的对象是下标族 $C_q(A)$。\par

\noindent\hangindent=2.2em\hangafter=1\textbf{p.210}\quad 首行把 $\pi_{n+3}(S^n)$ 排作 $\pi_{n+3}(S_n)$（$S$ 带下标 $n$），按 p.209 的记法应为上标。\par

\noindent\hangindent=2.2em\hangafter=1\textbf{p.211}\quad 例子编号重复：$\Omega^O_2$ 与 $\Omega^O_3$ 两条都标作 3.，后者按顺序应为 4.。\par

\noindent\hangindent=2.2em\hangafter=1\textbf{p.215}\quad 把 $G$-流形与配边的上标排作 $m$（$W^m_1$、$V^{m+1}$），按定义其维数应为 $k$；末式又印 $S^{n+k}$ 而上文用 $S^{m+k}$。\par

\noindent\hangindent=2.2em\hangafter=1\textbf{p.219}\quad 印作 $T(\CP^k)=1,\ j\ge0$，其中 $j$ 无定义；按上下文应为 $k\ge0$。\par

\noindent\hangindent=2.2em\hangafter=1\textbf{p.222}\quad $J$ 同态的像印作 $J\pi_{4k-1}(SO)\subset\pi_{4k-1}(S^n)$（$n>4k$），但此时右端为 $0$，与断言矛盾；按其本意应为稳定球面同伦群 $\pi_{n+4k-1}(S^n)$。\par

\noindent\hangindent=2.2em\hangafter=1\textbf{p.226}\quad 印作 「yields up the structures」，多一 「up」，系排印衍字。\par

\noindent\hangindent=2.2em\hangafter=1\textbf{p.227}\quad 两处把点写成 $s_0$，而同一处邻域记作 $U_x$；由 $U_x$ 判断该点应为 $x_0$。\par

\noindent\hangindent=2.2em\hangafter=1\textbf{p.232}\quad 印作 $\tau_j(P^{4k})=0$；按 Milnor 的构造，需要的是 $2k$ 维以下的同伦群为零（手术才能把 $P^{4k}$ 化为 $2k$-连通、其边界才是同伦球面），应为 $\pi_j(P^{4k})=0$。\par

\noindent\hangindent=2.2em\hangafter=1\textbf{p.234}\quad 印作 $bP^2\cong bP^6\cong bP^{14}\cong\Z/2$；但 $bP^2\subset\Theta^1=0$，且 $k=0,1,3$ 正是 Arf 不变量非平凡、Kervaire 不变量为 $1$ 的维数，按紧接的 Kervaire--Milnor 定理（$bP^{4k+2}\cong\Z/2$ 当且仅当 $\Phi=0$）应为 $0$。\par

\noindent\hangindent=2.2em\hangafter=1\textbf{p.243}\quad 式 (4.19) 记 $\pi_j(SPL/SO)=\Gamma^{j-1}$，而本段又说 $j$ 维骨架上的障碍属于 $H^j(M;\Gamma^j)$，两处 $\Gamma$ 的指标差 $1$；计算时以 (4.19) 的显式同构为准。\par

\noindent\hangindent=2.2em\hangafter=1\textbf{p.273}\quad 「Bernshtein」「Kirilov」是俄国学者姓名的转写异体（今多写作 Bernstein、Kirillov），并非原书排印错误；此条收录仅为便于读者按今通行拼法检索文献。\par

\noindent\hangindent=2.2em\hangafter=1\textbf{p.274}\quad 印作 「Calaby--Yau」，通行拼法为 Calabi--Yau；同段 「Salomon」通行拼法为 Salamon。\par

\noindent\hangindent=2.2em\hangafter=1\textbf{p.276}\quad 印作 $T^2\subset M^2$；$M$ 是纽结的补空间，为 3 维开流形，应为 $T^2\subset M^3$。\par

\noindent\hangindent=2.2em\hangafter=1\textbf{p.277}\quad 把因子印作大写的 $A^{\pm1}$，而同一段的 $Q_{D_0}$、$K_{D_n}$ 诸式皆用 $a$，应为 $a^{\pm1}$。\par

\noindent\hangindent=2.2em\hangafter=1\textbf{p.285}\quad 图 A.9 右半把柱体上底标为 $z=0$、下底标为 $z=1$，与正文「$n$ 辫的起点在底面 $z=0$、终点在顶面 $z=1$」正好相反；按正文，$z$ 增大应向上。\par

\noindent\hangindent=2.2em\hangafter=1\textbf{p.286}\quad 求和上下限写作 $j=0$ 至 $j=n$，而通项里用的是 $k$；按外幂的标准恒等式，求和指标应统一为 $k$。\par

\noindent\hangindent=2.2em\hangafter=1\textbf{p.287}\quad Hecke 代数关系印作 $\sigma_i-\sigma_{i-1}=z$；辫群生成元之间并无此关系，标准关系是 $\sigma_i-\sigma_i^{-1}=z$（同页 Birman--Wenzl 即作 $\sigma_i-\sigma_i^{-1}=z(v_i-1)$），系上标 $-1$ 误排成下标 $i-1$。\par

\noindent\hangindent=2.2em\hangafter=1\textbf{p.288}\quad 开尔文勋爵的姓印作 Thompson，通行拼法为 Thomson；同处 $J(t,t^{-1})$ 把 $t$ 与 $t^{-1}$ 并列成两个自变量，实为 $t$ 的 Laurent 多项式。\par

\noindent\hangindent=2.2em\hangafter=1\textbf{p.292}\quad Chern--Simons 作用量印成「局部形式」，式中下标 $i,j$ 未按 Einstein 约定缩并，$\partial_iA_i$ 亦非坐标不变的写法；标准形式为 $S(A)=\int_M\Tr\bigl(A\wedge dA+\frac23A\wedge A\wedge A\bigr)$。\par

\noindent\hangindent=2.2em\hangafter=1\textbf{p.296 两处}\quad (1) 把 Dijkgraaf 拼作「Dijgraaf」（p.296 末段与 p.297 亦如此）。(2) 本段写作 $V=S_1$，把上标 $1$ 误排成下标；前文一直写 $V=S^1$。\par

\noindent\hangindent=2.2em\hangafter=1\textbf{p.300}\quad Modern Geometry 英译本卷号印作「93 and 1104」；同一部书在 p.301（III.3）印作「93 and 104」，可证前者为「104」之误排。\par

\noindent\hangindent=2.2em\hangafter=1\textbf{p.301}\quad Hirzebruch 1956 年德文著作的丛书号印作「Ergeb. Math. g.」；同一部书在 p.303 印作「Ergeb. Math. 9」，可证「g.」为「9」之误排。\par

\noindent\hangindent=2.2em\hangafter=1\textbf{p.302}\quad 分组编号由 IV.2 直接跳到 IV.4，缺 IV.3。\par

\noindent\hangindent=2.2em\hangafter=1\textbf{p.303}\quad 「Characteristic Classs」多一个 s（第 II 册同一条目作 Classes），系排印之误。\par

\noindent\hangindent=2.2em\hangafter=1\textbf{p.304}\quad Dynin 一文的英译本年份印作 (1960)，而俄文原刊年份为 1966，应为 (1966)。\par

\noindent\hangindent=2.2em\hangafter=1\textbf{p.305 两处}\quad (1) Volodin--Fomenko 一文的刊名中「Voktorn」应为「Vektorn」。(2) Arnol'd 等书的出版者排作 Birkhauser（变音符脱字），应为 Birkhäuser。\par

\noindent\hangindent=2.2em\hangafter=1\textbf{p.310}\quad Kervaire 1960 年一文题名印作 「A manifolds which does not admit」，应为 「A manifold」；同条 Mather 一文的卷号年份印作 「29, No.\ 5 (167) (1972)」，「167」应为 1971。\par

\end{multicols}

\noindent 本表共收录 46 条，均经原书页图逐条核对。凡本表未列之处，或为原书无误，或为我们尚未察觉；正文与索引的排印，凡本表未列者，一律照原书排印。
\normalsize

% CHUNK erratatable DONE
